\section{Secure Lifecycle Plan}
\label{sec:lifecycle}

SecureBank~PT is developed in two-week sprints, and security activities
are part of each sprint and of the \ac{CI/CD} pipeline. The plan selects practices from NIST \ac{SSDF}~\cite{nist800218},
Microsoft \ac{SDL}~\cite{mssdl} and McGraw's touchpoints~\cite{mcgraw2006},
and uses the OWASP \ac{MASVS}~\cite{masvs} as the verification baseline for
the mobile application. Maturity is measured yearly with OWASP \ac{SAMM}~\cite{samm}, aiming at
level~2 in every practice. Table~\ref{tab:lifecycle} summarises the plan.

\subsection{Roles}
The \textbf{Security Architect (SA)} owns the threat model and design
reviews. The \textbf{Security Champion (SC)} is one developer per squad
who triages tool findings and runs code reviews, while \textbf{Developers
(DEV)} implement features and fix findings. The \textbf{QA/Security Tester
(QA)} runs dynamic tests and fuzzing and coordinates the external pentest. The \textbf{DevSecOps Engineer (DSO)}
manages the pipeline, signing keys, secrets and monitoring. The
\textbf{Product Owner (PO)} accepts security requirements and residual
risk. The \textbf{Compliance Officer / DPO (CO)} handles regulatory
mapping, and the \textbf{\ac{PSIRT}}, chaired by the bank CISO, handles
vulnerabilities and incidents after release.

\subsection{Training}
Before committing code, developers complete training on the OWASP Mobile
and \ac{API} Top~10, testers on \ac{MASVS}, and product owners on abuse
cases and MB~WAY fraud. It is repeated yearly and after major incidents.

\subsection{Requirements}
Requirements R1--R9 are written as testable user stories, for example ``an
access token older than 5~minutes is rejected''. Each function has at
least one abuse case~\cite{mcgraw2006}, such as enrolling a second device
with a stolen \ac{PIN}, replaying a signed transfer, or a coerced customer
sending money to a mule account. The CO maps them to \ac{PSD2},
\ac{GDPR} and \ac{CRA} Annex~I in a traceability matrix. The requirements
are then refined into measurable non-functional requirements, which set
acceptance thresholds rather than implementation choices and are verified
in the pipeline or before release (Table~\ref{tab:nfr}).

\begin{table}[!htbp]
\caption{Measurable non-functional requirements (S: security, A: availability, I: integrity)}
\label{tab:nfr}
\centering
\scriptsize
\begin{tabularx}{\columnwidth}{@{}L{0.75cm}L{4.15cm}Y@{}}
\toprule
\textbf{ID} & \textbf{Acceptance target} & \textbf{Traceability and evidence} \\
\midrule
NFR-S1 & 100\% of protected endpoints enforce authenticated and authorised access; zero unauthorised requests in negative tests. & R1, R3, R4; \ac{API} security and access-control reports \\
NFR-S2 & 100\% of sensitive operations require device-bound \ac{MFA} and a signature covering the operation data; replay and altered-data tests fail. & R2, R4; mobile/API tests and \ac{MASVS} results \\
NFR-S3 & No plaintext sensitive data or secrets in release artefacts; TLS~1.3 in transit and AES-256-GCM, or an approved equivalent, at rest. & R5, R8; MobSF, secret-scan and configuration reports \\
NFR-S4 & Every authentication attempt is logged with outcome and correlation ID; brute-force detection raises a \ac{SIEM} alert within 60 seconds. & R9; log-schema test and alert drill \\
NFR-A1 & At 1,000 concurrent sessions and 100 requests/s, 99.9\% of valid requests complete within 2 seconds, without mixing users or sessions. & R6; k6 load report and isolation assertions \\
NFR-A2 & A failed transfer leaves no partial state; recovery restores service within 30 minutes and loses no committed transaction. & R7; fault-injection, recovery and reconciliation reports \\
NFR-I1 & Each transfer is processed exactly once for a valid idempotency key; duplicate, replayed or tampered requests are rejected and audited. & R7, R9; \ac{API} replay/fuzz tests and audit-log review \\
NFR-I2 & A release is deployable only when its signed \ac{SBOM} and artefact signatures verify successfully on every node. & R8; \ac{CI/CD} gate and deployment verification record \\
\bottomrule
\end{tabularx}
\end{table}

\subsection{Design}
The SA applies \ac{STRIDE} to the \ac{DFD} of Figure~\ref{fig:context}
whenever a feature crosses a trust boundary~\cite{shostack2014}. Each
threat is assigned an owner and a mitigation, for example R2 and R4 for
spoofing, per-operation ECDSA signatures over the amount and payee for
tampering, and the signed audit log for repudiation. Design reviews use the
IEEE top ten design flaws~\cite{ieeecsd2014} and Table~\ref{tab:principles}
as checklists, and verify idempotency keys and exactly-once transfers. Only
vetted libraries and standard algorithms are used, and the update channel
is designed so that patches can be deployed quickly.

\subsection{Implementation}
Every merge request runs \ac{SAST}, MobSF, secret scanning and dependency
scanning, and a high or critical finding blocks the merge unless the PO
formally accepts the risk. Each build generates a CycloneDX \ac{SBOM}, so
the team knows which components run in each environment, which was the
inventory missing in the Equifax case~\cite{equifax2018}. Changes to
security-relevant code are reviewed by a second developer and by the SC.
Open findings per severity and time to fix are reported to the PO in each
sprint.

\subsection{Verification}
OWASP ZAP scans the \ac{API} every night, while MobSF and Frida check root
detection, certificate pinning and secure storage on the device. Jazzer
fuzzes the request parsers of the backend, and k6 load tests confirm R6 by
checking that the gateway rejects requests under overload instead of mixing
sessions. An authorised external team performs a penetration test before
each major release. Open findings go to a residual-risk list signed by
the PO.

\subsection{Release and Operations}
A go/no-go review checks the pipeline gates, the pentest status and the
\ac{SBOM} before each release. Artefacts are signed with keys in the
\ac{HSM}, and deployments are progressive and verified on every node, so a
partial release as in the Knight Capital case~\cite{knight2013} is
detected before reaching customers. The \ac{SIEM}
raises brute-force and fraud alerts (R9). The
\ac{PSIRT} publishes a vulnerability disclosure policy, releases critical
fixes within 72~hours as security-only updates, and rehearses the \ac{IRP}
twice a year, including the \ac{CRA} and \ac{DORA} reporting deadlines.

\begin{table}[!htbp]
\caption{Secure lifecycle plan: activities, owners, tools and evidence per phase}
\label{tab:lifecycle}
\centering
\scriptsize
\begin{tabularx}{\columnwidth}{@{}L{1.55cm}Y L{0.95cm}Y@{}}
\toprule
\textbf{Phase} & \textbf{Activities} & \textbf{Owner} & \textbf{Tools and evidence} \\
\midrule
Training & Role-based secure coding; mobile and \ac{API} pitfalls; social-engineering awareness; yearly refresh & SC, CO & OWASP MASTG labs, internal LMS; training records, quiz scores \\
Requirements & R1--R9 as user stories; abuse cases; compliance mapping (\ac{PSD2}, \ac{GDPR}, \ac{CRA}) & PO, SA, CO & Jira, \ac{MASVS} L2 and R; requirements catalogue, traceability matrix \\
Design & \ac{STRIDE} on the \ac{DFD}; architecture review; key management design & SA, SC & OWASP Threat Dragon; threat model, review minutes, crypto inventory \\
Implementation & Secure coding standard; peer review; \ac{SAST}, dependency and secret scanning; \ac{SBOM} & DEV, SC, DSO & Semgrep, MobSF, Dependency-Track, Gitleaks, CycloneDX; review logs, signed \ac{SBOM} \\
Verification & \ac{DAST} of the \ac{API}; mobile dynamic tests; fuzzing; load tests; external pentest & QA, external & OWASP ZAP, Frida, Jazzer, k6; test and pentest reports, residual-risk list \\
Release \& Operations & Go/no-go review; artefact signing; \ac{SIEM} monitoring; disclosure and patching; \ac{IRP} & PO, DSO, \ac{PSIRT} & Sigstore, Wazuh, security.txt; go/no-go record, advisories, patch metrics \\
\bottomrule
\end{tabularx}
\end{table}